% Ad Familiares 1.2 --- headnote
% ad-familiares-01-02, 17 January 56 BC, written at Rome

\textit{Cicero to Lentulus Spinther, written from Rome before
dawn on the sixteenth day before the Kalends of February
(17 January) 56 BC, four days after Fam. 1.1. The dispatch
follows the Senate from the Ides into the next session: the
Bibulus motion (three legates), the Hortensius motion (Lentulus
without an army), the Volcacius motion (Pompey). The procedural
move that broke the day was the demand to divide Bibulus's
motion --- the religious half (the Sibyl: no army) was assented
to as inevitable, the three-legates half voted down. With
Hortensius's motion next on the table, the tribune Lupus,
having previously moved Pompey, claimed precedence in calling
the division. The consuls let the day run out without forcing
the issue; they wanted Bibulus's opinion to stand. Cicero's
private take on Pompey is the central note --- ``when I hear
him alone, I plainly clear him of all suspicion of grasping;
but when I see his intimates, I see clearly that this whole
business has long been corrupted by certain men, the king
himself and his counsellors not unwilling.'' The closing
\textsection 4 reports what Cicero counts as the day's gain:
a senatus auctoritas (intercession by Cato and Caninius
notwithstanding) which forecloses any popular legislation
without breaking auspices, laws, or peace. The letter is the
clearest single statement of how the senatorial centre is
holding by procedure --- religion, the auspices, the proxies
of Cato and Hortensius --- against a Pompey who in private
talks like a friend and in his entourage acts like a rival.}
